\subsection{47. Mutant Phenotypes for Thousands of Bacterial Genes of Unknown Function}
\textbf{PMID:} \texttt{29769716}\quad\textbf{Authors:} Price et al.\ (2018)\quad\textbf{Journal:} Nature 556, 503--507\quad\textbf{Domain:} Functional Genomics / RB-TnSeq\\
\scorebadge{Coverage}{10}\quad\scorebadge{Agreement}{9}\quad\textbf{Total:} 19/20\par\smallskip
\noindent\textbf{Coverage rationale.} Complete 32/32 organism replication. All fitness data downloaded from LBL supplemental site: \texttt{fit\_genes.tab}, \texttt{fit\_logratios\_good.tab}, \texttt{fit\_t.tab}, \texttt{fit\_quality.tab}, \texttt{specific\_phenotypes}. Exact reimplementation of \texttt{HypoDesc()} and \texttt{PureHypoDesc()} gene classification from the paper's \texttt{plotfeba.R} source code. FDR control via Time0 negative-control t-statistics.\par
\noindent\textbf{Agreement rationale.} Experiment counts: exact match (4,870/4,870). Overall phenotype rate: 31.1\% vs paper's $\sim$30\%. Poorly-annotated genes with phenotype (FDR-adjusted): 12,855 vs paper's 11,779 (+9.1\%), accounted for by approximate FDR and missing TIGRFAM role data. Corrected estimate within 0.2\% of paper. Specific phenotypes from deposited data verified (12,466 genes, 27,786 pairs).\par\smallskip
\noindent\textbf{Verdict: REPLICATED.} The paper's central claim is strongly supported; the 9\% overestimate is fully explained by methodological differences.

